% ================= 2. SYSTEM REQUIREMENTS =================
\section{System Requirements and Verification}
\label{sec:requirements}

Based on the objectives established in Section~\ref{sec:objectives}, a set of technical requirements has been developed to ensure the final design aligns with overall project goals. While a portion of these requirements flows directly from top-level objectives, others were derived during preliminary sizing to accommodate constraints identified during modeling and trade studies.

The complete requirement set is summarized in Table~\ref{tab:requirements}. To parse this matrix, Table~\ref{tab:req_key} defines the verification codes used to demonstrate compliance, while full testing procedures, acceptance criteria, and evidence are detailed in Appendix~\ref{app:verification}.
\begin{table}[htbp]
\centering
\caption[Verification methods]{Verification methods used in
Table~\ref{tab:requirements}.}
\label{tab:req_key}
\small
\setlength{\tabcolsep}{5pt}
\begin{tabular}{@{} c l L{9.0cm} @{}}
\toprule
Code & Name & Meaning \\
\midrule
A & Analysis & Closure by calculation or simulation against a stated model. \\
I & Inspection & Confirmation against drawings, a parts list, or a measurement of the built article. \\
T & Test & Instrumented measurement against a numerical criterion. \\
D & Demonstration & Observed execution of a behaviour, without instrumented measurement. \\
\bottomrule
\end{tabular}
\end{table}

\begin{small}
\setlength{\tabcolsep}{4pt}
\begin{longtable}{@{} l L{7.5cm} L{2.5cm} c @{}}
\caption[System requirements]{System requirements, with parent objective and
verification method. (D) marks a derived requirement. Methods: A analysis,
I inspection, T test, D demonstration.}
\label{tab:requirements}\\
\toprule
ID & Requirement & Parent objective & V \\
\midrule
\endfirsthead
\multicolumn{4}{c}{\footnotesize\itshape \tablename~\thetable\ (continued)}\\
\toprule
ID & Requirement & Parent objective & V \\
\midrule
\endhead
\midrule
\multicolumn{4}{r}{\footnotesize\itshape Continued on next page}\\
\endfoot
\bottomrule
\endlastfoot

\multicolumn{4}{@{}l}{\textbf{Functional}}\\
\midrule
FUN-01 & The cube shall maintain balance about an edge equilibrium using only internal reaction-wheel actuation. & Edge balance & T \\
FUN-02 & The cube shall maintain balance about a corner equilibrium under three-axis feedback. & Corner balance & T \\
FUN-03 & The cube shall execute a commanded slew between two specified attitudes and re-capture balance at the target. & Controlled rotation & T \\
FUN-04 & The cube shall reject an impulse disturbance applied to the body and return to the balancing equilibrium. & Edge, corner balance & T \\
FUN-05 & The system shall estimate body attitude and rate from onboard sensing alone, with no external reference. & All & T \\
FUN-06 & The system shall bound stored wheel momentum by commanding a desaturation torque above a defined wheel-speed threshold. & All & T \\
FUN-07 & The system shall sequence operating modes through a supervisory state machine with hysteresis on every threshold. & All & D \\
FUN-08 & The system shall provide a disarm transition, reachable from every operating mode, that de-energises the drivers. & All & D \\
FUN-09 & On loss of the balancing equilibrium the system shall enter a safe state with the wheels commanded down, rather than continuing to actuate. & All & D \\
FUN-10 & The cube shall transfer stored wheel momentum through a braking event to tip from a face rest onto an edge, and capture the resulting attitude into balance. & Jump-up & T \\
FUN-11 & The cube shall execute consecutive jump-up manoeuvres, re-stabilising after each landing. & Walking & T \\
FUN-12 & The system shall record telemetry without perturbing control-loop timing. & All & D \\
\midrule

\multicolumn{4}{@{}l}{\textbf{Performance}}\\
\midrule
PER-01 & The system shall maintain edge balance for \qty{60}{\second} or more unassisted, with no operator contact and no evident wheel-speed drift, on the 1-DoF rig and on the integrated cube. & Edge balance & T \\
PER-02 & The cube shall maintain corner balance for \qty{60}{\second} or more unassisted, with no evident wheel-speed drift. & Corner balance & T \\
PER-03 & The 1-DoF rig shall recover to the balancing equilibrium from an initial tilt of \qty{12}{\degree} or more, released from rest, re-establishing PER-01. & Edge balance & T \\
PER-04 & The cube shall recover to the balancing equilibrium from an initial tilt of \qty{5}{\degree} or more in any direction, re-establishing PER-01 or PER-02 as applicable. & Edge, corner balance & T \\
PER-05 & The control and estimation loop shall execute at \qty{1}{\kilo\hertz} with a worst-case cycle time not exceeding \qty{1}{\milli\second}. & All & T \\
PER-06 (D) & Each reaction wheel shall provide an inertia of at least \qty{6.2352e-4}{\kilogram\metre\squared} about its spin axis. & Jump-up & A, I \\
PER-07 & The drive shall sustain the design wheel speed of \qty{4000}{\rpm} at end-of-discharge pack voltage. & Jump-up & A, T \\
PER-08 & The design wheel speed shall not exceed \qty{80}{\percent} of the achievable maximum wheel speed at end-of-discharge pack voltage. & Jump-up & A \\
PER-09 & Commanded motor current shall remain within the driver's \qty{9}{\ampere} continuous rating throughout balancing. & All & T \\
PER-10 & The brake shall arrest a wheel from the design wheel speed within \qty{100}{\milli\second}. & Jump-up & T \\
PER-11 & Telemetry shall record the state estimate, the control command and the wheel speeds at \qty{50}{\hertz} or faster. & All & T \\
\midrule

\multicolumn{4}{@{}l}{\textbf{Design constraint}}\\
\midrule
CON-01 & The cube outer edge length shall be \qty{150}{\milli\metre}. & All & I \\
CON-02 & Total cube mass shall not exceed \qty{1.85}{\kilo\gram}. & Jump-up & I \\
CON-03 & The system shall be actuated by exactly three mutually orthogonal reaction wheels. & All & I \\
CON-04 & The cube shall operate from a single onboard 6S lithium-polymer pack, with no tether or external power during a demonstration. & All & I, D \\
CON-05 & Driver communication shall use a single linear CAN-FD chain, terminated at exactly the two physical ends. & All & I \\
CON-06 & The \qty{5}{\volt} logic rail shall carry its coincident worst-case load within the regulator's derated rating. & All & A, T \\
CON-07 & The electronics shall fit within the frozen board outline, mounting pattern and keep-out heights of the electronics-mount interface. & All & I \\
\end{longtable}
\end{small}
